\subsection{Causal Discovery}
\totalscore{7}

Let $M$ be a simple SCM with observed endogenous variables $O = \{1,2,3\} \subseteq V$ and (observed) exogenous input variables $J = \{4\}$, with observable graph $G^{O\cup J}(M)$.
Suppose that $M$ is $\sigma$-faithful with respect to $O$.
Assume further that $X_2$ is not a cause of $X_1$ according to $M$.

You observe the following conditional (in)dependencies in data sampled from the Markov kernel $P_M(X_O \given \doit(X_4))$:
  $$X_1 \nIndep X_2 \given X_4 \quad \land \quad X_2 \nIndep X_3 \given X_4 \quad \land \quad X_1 \Indep X_3 \given X_2,X_4$$

\begin{enumerate}[label=\alph*)]
  \item Draw all possible observational graphs $G^{O\cup J}(M)$ that are compatible with the assumptions and observations, 
    and explain why no other graph is possible.
    You may use dashed edges to indicate an edge that could be present or absent.\score{5}
    \answer{Any edge between $X_4$ and another variable must be out of $X_4$. \point{(1 pt)}\\
The assumption that $X_2$ is not a cause of $X_1$ implies that there is no directed edge $X_2 \tuh X_1$ in the graph. \point{(1 pt)} \\
    If there were an edge between $X_1$ and $X_3$, $1 \nPerp 3 \given \{2, 4\}$ and hence, by faithfulness, $X_1 \nIndep X_3 \given X_2, X_4$ which contradicts the observed independence. \point{(1pt)}\\
    Also, since $X_1 \nIndep X_2 \given X_4$, by the Markov property, $1 \nPerp 2 \given 4$ which implies that $X_1$ and $X_2$ must be adjacent. Similarly, $X_2$ and $X_3$ must be adjacent. \point{(1pt)}\\
$X_2$ cannot be a collider on any walk between $X_1$ and $X_3$ (faithfulness).
    Since the only possible edges between $X_1$ and $X_2$ are $X_1 \tuh X_2$ and $X_1 \huh X_2$ (both of which have an arrowhead at $X_2$), this means that there must be a directed edge $X_2 \tuh X_3$, but there cannot be a  bidirected edge $X_2 \huh X_3$ or directed edge $X_2 \hut X_3$. \point{(1pt)}\\
    Thus, the only possible graphs are the following:\\[\baselineskip]
\centerline{\begin{tikzpicture}
    \begin{scope}
      \node[ndout] (X1) at (-3,0) {$X_1$};
      \node[ndout] (X2) at (-1.5,0) {$X_2$};
      \node[ndout] (X3) at (0,0) {$X_3$};
      \node[ndint] (X4) at (-1.5,1.5) {$X_4$};
      \draw[arout] (X1) edge (X2);
      \draw[arout] (X2) edge (X3);
      \draw[arint,dashed] (X4) edge (X1);
      \draw[arint,dashed] (X4) edge (X2);
      \draw[arint,dashed] (X4) edge (X3);
    \end{scope}
    \begin{scope}[xshift=4.5cm]
      \node[ndout] (X1) at (-3,0) {$X_1$};
      \node[ndout] (X2) at (-1.5,0) {$X_2$};
      \node[ndout] (X3) at (0,0) {$X_3$};
      \node[ndint] (X4) at (-1.5,1.5) {$X_4$};
      \draw[arlat,bend left] (X1) edge (X2);
      \draw[arout] (X2) edge (X3);
      \draw[arint,dashed] (X4) edge (X1);
      \draw[arint,dashed] (X4) edge (X2);
      \draw[arint,dashed] (X4) edge (X3);
    \end{scope}
    \begin{scope}[xshift=9cm]
      \node[ndout] (X1) at (-3,0) {$X_1$};
      \node[ndout] (X2) at (-1.5,0) {$X_2$};
      \node[ndout] (X3) at (0,0) {$X_3$};
      \node[ndint] (X4) at (-1.5,1.5) {$X_4$};
      \draw[arout] (X1) edge (X2);
      \draw[arlat,bend left] (X1) edge (X2);
      \draw[arout] (X2) edge (X3);
      \draw[arint,dashed] (X4) edge (X1);
      \draw[arint,dashed] (X4) edge (X2);
      \draw[arint,dashed] (X4) edge (X3);
    \end{scope}
    \end{tikzpicture}}
    }
%  \item Derive the most informative DPAG (i.e., with as few circles as possible) that represents all of the possible DMGs you found in a).
  \item Can the causal effect of $X_2$ and $X_4$ on $X_3$, i.e., the quantity
    $$P_M(X_3 \given \doit(X_4), \doit(X_2))$$
    be identified from the Markov kernel $P_M(X_O \given \doit(X_4))$?
    If not, explain why not. If yes, provide an expression for the causal effect in terms of the observed Markov kernel and give a corresponding derivation, explaining why your identity must hold.\score{2}
  \answer{Yes, it can be identified. 
    The second rule of the causal do-calculus applied to $G^{O\cup J}(M)$ yields:
      \begin{equation*}
        P_{M}\big(X_3 \given \doit(X_4), \doit(X_2)\big) \stackrel{a.s.}{=} P_{M}(X_3 \given \doit(X_4), X_2),
      \end{equation*}
      because
      $$3 \Perp^\sigma_{(G^{O\cup J}(M))_{\doit(I_2)}} I_2 \given \{2, 4\}.$$
  }
  \points{1 point for the right identity, 1 point for the derivation.}
\end{enumerate}
